% ---------------------------------------------------------------------------
% Draft subsection: the end-to-end pipeline.
% Written 2026-07-26 in response to Nicolo's reply of the same date
% ("Can we draft a subsection on the end-to-end pipeline with all these
% comments? Would the figure I proposed above fit in it?").
%
% Numbers marked [RERUN] are placeholders pending the full-scale RHUL job
% slurm/nicolo_20260726.sbatch.  Every other number in this file has been
% measured on this tree; provenance is given in the comment above it.
% Do NOT paste [RERUN] values into Overleaf until the job has landed.
% ---------------------------------------------------------------------------

\subsection{The end-to-end pipeline}
\label{sec:end-to-end}

The method chains five deterministic stages: an empirical prior~$H$ fitted to
catalogued systems, a generative noise model, a synthetic corpus, an amortised
regressor~$\psi$, and a conformal calibration layer.  This subsection states
what each stage consumes, and --- more importantly --- what it is \emph{not}
allowed to see, since the validity of the final intervals rests entirely on
that bookkeeping.

\paragraph{What the model is trained on.}
$\psi$ is trained exclusively on synthetic curves, with the data-generating
parameters~$\bar\theta$ as labels.  No tabulated parameter is ever a training
target.  Catalogued values enter in exactly four places, and in the first three
only through the \emph{training} split of the real corpus: (i)~fitting the
empirical histograms~$H$ over $\log_{10}P$, $\log_{10}K$ and~$e$; (ii)~the
paired (cadence,~$\sigma$) bootstrap that supplies observation times; (iii)~the
residual set on which the SVGP noise model is fitted; and (iv)~at test time, as
the initialiser for the gradient-descent surrogate~$\theta^\ast$ on real
curves.  Validation and test systems contribute to none of (i)--(iii), so the
distribution the synthetic corpus imitates is shaped only by systems we are
never scored on.

% Measured 2026-07-26 on this tree, seed 42, 200 epochs, feature set 74:
%   real-split all   R2 = 0.847   (357 series)
%   real-split train R2 = 0.856   (242 series)
%   real-split test  R2 = 0.804   ( 57 series)
Synthetic-to-real transfer is correspondingly honest: regressing
$\log_{10}K$ on the held-out test split gives $R^2 = 0.804$, against $0.856$ on
the training split whose parameters shaped~$H$.  We quote the test-split figure
throughout; the gap is the price of the distribution-level leakage that
fitting~$H$ necessarily incurs.

\paragraph{Which parameters are reported.}
% Nicolo 2026-07-26: "Let's skip showing and using omega as a target."
We report prediction sets for $P$, $K$ and~$e$ only.  The argument of
periastron~$\omega$ is not identifiable at the signal-to-noise ratio of this
corpus: its conformal half-width exceeds~$\pi$, so the nominal interval wraps
the entire circle and communicates nothing.  Rather than tabulate a number that
reads like a measurement, we state the non-identifiability and drop~$\omega$
from the reported coordinates, which also reduces the Bonferroni divisor from
four to three.  Note that $\omega$ is \emph{not} removed from the model: it is
still optimised, it still enters the decoder, and fixing it --- for instance at
a catalogue mean --- would be actively harmful, because the eccentricity is
parameterised through $(h,k) = (e\cos\omega, e\sin\omega)$ and the periastron
epoch~$T_{\mathrm{peri}}$ is refitted analytically against it.  Constraining an
unidentifiable angle would bias the two quantities we do report.

\paragraph{Which observations are usable.}
% Measured 2026-07-26: min_nights >= 5 drops 6 of 57 test series, leaving 51.
% The label-aware alternative (baseline >= one tabulated period) drops the same
% six.  Corpus-wide: 343 of 357 series survive.
The corpus stores one row per RV time series, and a small number of these are
transit-night sequences --- high-cadence Rossiter--McLaughlin runs covering a
single night.  HD~17156 contributes one series of 13 points spanning
$0.102$~days against a $21.2$~day period; such a series cannot constrain an
orbit at all.  We therefore require a series to span at least five distinct
nights.  The criterion is deliberately label-free, depending only on the
observing pattern and never on the tabulated period, so that test-set selection
does not consult the values we are later scored against; a label-aware
alternative (requiring the baseline to cover at least one tabulated period)
removes exactly the same six series from the test split.  After filtering,
$51$ of $57$ test series remain.

% 51 filtered test series span 34 distinct hosts; HD 179949 contributes 4.
These $51$ series come from only $34$ distinct stars, since a star observed by
several instruments contributes several series that share one true~$\theta$.
We therefore report coverage both at series level and with a single
best-sampled series per star, and quote the latter as the headline: it is the
variant whose test points are genuinely independent.

\paragraph{Calibration and what coverage is measured against.}
Calibration uses split conformal prediction on synthetic curves, with the
conformity score $s_c = |\psi_c(y) - \theta^\ast_c|$ against the
gradient-descent surrogate~$\theta^\ast$, reweighted by the likelihood ratio
$p(y_{\mathrm{real}})/p(y_{\mathrm{fake}})$ to correct the covariate shift
(Tibshirani et al., 2019).  On real systems, coverage is then evaluated against
the \emph{catalogued} parameters --- not against~$\theta^\ast$, and not against
any quantity our own optimiser produced.  This is the stronger of the two
possible claims: the reported intervals are checked against the values an
astronomer would look up, while $\theta^\ast$ appears only inside the synthetic
calibration step, where the data-generating~$\bar\theta$ is known.

\paragraph{Score normalisation.}
% Nicolo 2026-07-26: the re-encoding form avoids fitting any extra model.
Equation~(\ref{eq:delta}) is normalised by a per-system uncertainty proxy built
from two re-encoding residuals, neither of which requires training an
additional model.  The first, $\delta_c(y) = |\psi_c(h(\psi(y))) - \psi_c(y)|$,
passes the noiseless reconstruction back through the encoder and measures how
stably $\psi$ recovers its own prediction.  The second,
$\delta_y(y) = \|y - h(\psi(y))\|$, is the reconstruction residual in data
space.  Both are pointwise functions of a single curve, because $h$ and $\psi$
are deterministic.  They are combined convexly,
$w\,\delta_c + (1-w)\,\delta_y$, with $w$ and the regulariser~$\gamma$ selected
together on a tuning set by minimising interval width.
% [RERUN] The smoke-scale run (n_cal=40) selects w = 1.00, i.e. delta_y adds
% nothing once delta_c is present.  Confirm at n_cal = 400 before asserting.
The selected weight is reported alongside the intervals.

\paragraph{Finite-sample validity without curvature constants.}
% Nicolo 2026-07-26: "what can we say about the test sample (not used to
% compute the max)?"
Theorem~\ref{thm:main} carries a constant~$\Delta_c$ bounding the gap between
the data-generating parameter and the surrogate label.  Estimating it as a
maximum over a tuning set invites the objection that a fresh test point may
exceed it.  The objection is answered by treating $\Delta_c$ as a conformal
quantile rather than a plug-in constant: if the tuning gaps and the test gap
are exchangeable, the $k$-th order statistic with
$k = \lceil (1-\beta)(n+1) \rceil$ satisfies
$\Pr(|\bar\theta_c - \theta^\ast_c| > \Delta_c) \le \beta$ out of sample.  The
maximum is the case $k = n$, i.e.\ $\beta = 1/(n+1)$ --- so it is already
valid, at a cost of one part in $n+1$ of the error budget, under one per cent
at our tuning-set size.  Spending $\beta$ on the gap and $\alpha - \beta$ on
the main quantile yields a finite-sample statement that invokes no curvature
assumption at all.  On real data the same likelihood-ratio reweighting must be
applied to the gap, since exchangeability holds only after the shift is
corrected.

This matters because the curvature route is vacuous at the measured constants.
% kappa in the four physical coordinates ~3.6e3; the earlier ~1e6 figure was an
% artefact of differentiating in the redundant 5-vector.
With $\varepsilon$ on the scale of the loss itself, $\sqrt{C_H \varepsilon}$
exceeds the entire support of~$e$, so adding the corresponding correction would
inflate every interval to the whole parameter space.  We therefore report the
empirical constants for transparency and calibrate the intervals
conformally.  The theoretical bound establishes that the method is correct in
an idealised regime; it is the conformal construction, not the bound, that
makes the intervals we report valid.

\paragraph{Assumption on the loss geometry.}
% Measured 2026-07-26, 25 prior draws, 200 GD steps, physical coordinates:
%   PKew: PD fraction 0.64, median lambda_min +2.15, median kappa 2.26e4
%   PKe : PD fraction 0.76, median lambda_min +20.3, median kappa 6.38e3
% [RERUN] the 1000/4000-step sweep is in slurm/nicolo_20260726.sbatch.
Assumption~\ref{ass:hessian} asks that $\theta^\ast$ be a strict local minimum
of~(\ref{eq:objective}).  Measured at $\theta^\ast$ in the physical
coordinates, the Hessian is positive definite for $64\%$ of prior draws when
$\omega$ is included and $76\%$ when it is not; the median smallest eigenvalue
rises from $+2.1$ to $+20.3$ and the median condition number falls from
$2.3\times10^4$ to $6.4\times10^3$.  The assumption is thus substantially, but
not entirely, an artefact of the unidentifiable angle: descending into a nearly
flat direction in~$\omega$ leaves the optimiser at a point that is not a strict
minimum.  We report the residual failure rate rather than assuming it away.

One caveat on measurement: the five-vector
$(\log_{10}P, \log_{10}K, e, \cos\omega, \sin\omega)$ is a redundant encoding,
and the decoder sees~$\omega$ only through $\operatorname{atan2}$.  A Hessian
formed in those coordinates is singular by construction, with an exact null
direction along the $(\cos\omega, \sin\omega)$ radius.  All curvature
quantities above are computed in the physical coordinates instead.

\paragraph{Cost.}
% Measured 2026-07-26, CPU, 57 real test series, 256 observations each:
%   featurise 6.4 ms/system; psi forward 1.0 us/system; GD theta* 4.0 s/system
%   at 200 steps.  Training psi: ~14 s for 200 epochs on 10k rows.
The argument for amortisation is a cost argument, so it is worth stating in
absolute terms.  Inference on one real system costs $6.4$~ms, almost all of it
feature extraction: the forward pass through~$\psi$ itself is $1$~$\mu$s, and
training~$\psi$ takes $14$~seconds once for the whole corpus.  Fitting a single
system by gradient descent --- our own $\theta^\ast$, at 200 steps --- costs
$4.0$~s, roughly $600\times$ more, and that is an optimiser returning a point
estimate, not a sampler returning a posterior.  A Markov chain requires several
orders of magnitude more Keplerian integrations again.  This is the trade the
paper is making: the catalogued values are better fits, because they were
obtained by per-system inference on these exact data, and we are comparing a
single amortised forward pass against that.

\paragraph{Reconstruction quality, and its ceiling.}
% figures/paper/rv_mse_scatter.csv, 57 held-out planets, label-free scoring
% with t_peri and gamma refitted for both parameter sets.
Scored label-free --- each parameter set judged only on how well it
reconstructs the observed curve, with $T_{\mathrm{peri}}$ and~$\gamma$ refitted
for both so that neither is handicapped --- the median reconstruction MSE is
$0.113$ for the catalogued values and $1.16$ for ours, against a floor of
$0.078$ attained by gradient descent initialised at the catalogued values.
Three of $57$ planets fall on or below the diagonal.  We do not present this as
a win.  The catalogue sits within roughly $30\%$ of the achievable floor
because those values were fitted to these very data by the discovery teams;
the honest reading is that a single forward pass recovers orbits within an
order of magnitude of bespoke per-system inference, at the cost quoted above,
and returns a calibrated region rather than a point.

\paragraph{What the intervals buy.}
% [RERUN] all four numbers below pending the full-scale job.
Two comparisons make the case.  First, coverage: the reported regions attain
close to nominal coverage on real systems, whereas standard split conformal
calibrated on~$\bar\theta$ attains $1.000$ --- valid, but far too wide to be
useful.  Second, containment against the catalogue: our $\alpha = 0.1$ region
contains the catalogued value for the large majority of planets, while the
catalogue's own $90\%$ interval contains our prediction far less often.  The
asymmetry is the point.  It shows that the intervals are absorbing the
synthetic-to-real distribution shift rather than merely being wide, which is
precisely what the reweighted calibration is designed to do.

% Nicolo asked whether the centred 2-D figure fits in this subsection: yes,
% here, illustrating the containment claim above.
Figure~\ref{fig:regions} shows this directly.  Because the regions differ by
some two orders of magnitude in~$P$ between the two methods, each planet is
centred on its own catalogued value --- axes are $\log_{10}(x/x_{\mathrm{tab}})$
for $P$ and~$K$ --- so that all ten systems share one frame.  The conformal
regions appear as legible outlines and the catalogue regions collapse towards
the origin.  We show the $(K, e)$ panel rather than $(K, \omega)$, for the
identifiability reason given above.
